\documentclass[aps,prd,onecolumn,amsmath,amssymb,longbibliography,floatfix,nofootinbib]{revtex4-2}

\usepackage{hyperref}
\usepackage{url}
\usepackage{booktabs}
\usepackage{longtable}
\def\UrlBreaks{\do\/\do\_\do\.\do\:}
\newcommand{\cell}[2]{\parbox[t]{#1\linewidth}{\raggedright #2\strut}}
% Generated by papers/two_center_born_infeld/generate_numbers.py; do not edit.
% ref acorr28 = origin/research/acorr-newton-g-c28 at 20cb3828b3550103544a0ae5dee6af2df28b6bb6
% ref bond = origin/substrate-bond-generator-20261101 at 367287a67291015d4b68f7cdfe2eb1119aa39d26
% ref c27 = origin/review/order17-divisor-c27-20261101 at cdd0f70ffa18e5599134127805c555ba2fec76cb
% ref dev = origin/dev at 830cc5170df58c182ce54fe7a8c7ca0cdb2e3842
% ref ion = origin/research/ionization-generator-period3-20261101 at 40cd2c13f73b5747302d2d70dd81dfabf08dd452
% ref ode = origin/math/rtail-frozen-ode-1-10 at 3b8c2653996dda0629c544328b48abc5ef0ca38b
\newcommand{\RefAcorrTwoEight}{\texttt{20cb3828b355}}
\newcommand{\RefBond}{\texttt{367287a67291}}
\newcommand{\RefCTwoSeven}{\texttt{cdd0f70ffa18}}
\newcommand{\RefDev}{\texttt{830cc5170df5}}
\newcommand{\RefIon}{\texttt{40cd2c13f73b}}
\newcommand{\RefOde}{\texttt{3b8c2653996d}}
\newcommand{\AcorrDiagVal}{-0.046663253999494962615}
\newcommand{\AcorrMagVal}{0.0466632539995}
\newcommand{\ApparentDegVal}{39}
\newcommand{\BondCoreRadiusVal}{\ensuremath{3.4831\times10^{-15}}}
\newcommand{\BondDiffMaxVal}{\ensuremath{2.1\times10^{-53}}}
\newcommand{\BondDiffMinVal}{\ensuremath{1.6\times10^{-85}}}
\newcommand{\BondFenceRows}{1/100 & 1.018981954 & 0.0101898 \\
1/4 & 0.9022344963 & 0.225559 \\
1 & 0.5886190197 & 0.588619 \\
4 & 0.1449942061 & 0.579977 \\
100 & 0.005027163156 & 0.502716 \\}
\newcommand{\BondHalfVal}{3.32720691976981241200687458504}
\newcommand{\BondImagVal}{\ensuremath{5.82\times10^{-86}}}
\newcommand{\BondScaleDVal}{28709.937}
\newcommand{\BondScaleRelVal}{\ensuremath{3.92\times10^{-14}}}
\newcommand{\BondSeriesDpsVal}{100}
\newcommand{\BondSeriesRows}{1.25 & 19 & $\ensuremath{1.63\times10^{-85}}$ \\
1.5 & 24 & $\ensuremath{4.35\times10^{-81}}$ \\
2 & 28 & $\ensuremath{1.2\times10^{-75}}$ \\
3 & 55 & $\ensuremath{3.59\times10^{-63}}$ \\
4 & 59 & $\ensuremath{2.07\times10^{-53}}$ \\}
\newcommand{\ContactBitsVal}{160}
\newcommand{\ContactTotalVal}{3.496076739056160}
\newcommand{\ContactWidthHundredthVal}{0.000227}
\newcommand{\CuspTailSumVal}{11.23}
\newcommand{\CzeroVal}{1.854074677301372}
\newcommand{\DivergenceColsVal}{19}
\newcommand{\DoneLowerVal}{0.63726946}
\newcommand{\DoneRelGapVal}{1.392}
\newcommand{\DoneUpperVal}{0.64626656}
\newcommand{\EndpointRatioVal}{0.5522847498307934}
\newcommand{\EnergyDegMaxVal}{192}
\newcommand{\EnergyDegMinVal}{166}
\newcommand{\EnergyOrderVal}{18}
\newcommand{\EpsTwoVal}{-0.15766642946915094105}
\newcommand{\EpsTwoZeroVal}{-0.12533}
\newcommand{\EtaTwoVal}{-0.111003175469655978435}
\newcommand{\ExactMinusScreenedVal}{0.0600102}
\newcommand{\FrozenCfourVal}{0.9270373386506860}
\newcommand{\FrozenMinusUpperVal}{0.05016}
\newcommand{\FrozenVoneVal}{0.696429977974}
\newcommand{\HDegMaxVal}{56}
\newcommand{\HDegMinVal}{40}
\newcommand{\HigherlFracVal}{69.3}
\newcommand{\IndInfEighteen}{\{-3,-1,-1/3,-1/3,0,1/3,1,1,5/3,2,7/3,7/3,4,5,6,7,9,9\}}
\newcommand{\IndInfSeventeen}{\{-3,-1,-1/3,-1/3,0,1/3,1,5/3,2,7/3,7/3,4,5,6,7,9,9\}}
\newcommand{\IndZeroEighteen}{\{-1,0,1,1,1,7/3,3,11/3,5,5,6,19/3,23/3,9,9,13,13,17\}}
\newcommand{\IndZeroSeventeen}{\{-1,0,1,1,1,7/3,11/3,5,5,6,19/3,23/3,9,9,13,13,17\}}
\newcommand{\IvpErrVal}{\ensuremath{1.07\times10^{-81}}}
\newcommand{\IvpOracleVal}{\ensuremath{3.99\times10^{-43}}}
\newcommand{\IvpTargetVal}{1.005}
\newcommand{\IvpThreshVal}{\ensuremath{1.0\times10^{-40}}}
\newcommand{\IvpWrongVal}{\ensuremath{1.0\times10^{-8}}}
\newcommand{\KVal}{1.236049784867581}
\newcommand{\KappaChangeVal}{\ensuremath{1.48\times10^{-68}}}
\newcommand{\KappaHalfGapVal}{0.160303}
\newcommand{\KappaSuperVal}{0.339697095903340977473156586990113338}
\newcommand{\MasterMaxErrVal}{\ensuremath{8.01\times10^{-15}}}
\newcommand{\NewFactorDegVal}{35}
\newcommand{\OldFactorDegVal}{39}
\newcommand{\OrderSeventeenVal}{17}
\newcommand{\ResidualRows}{1 & 96 & $\ensuremath{4.51\times10^{-53}}$ & $\ensuremath{2.03\times10^{-6}}$ \\
1.25 & 96 & $\ensuremath{2.63\times10^{-70}}$ & $\ensuremath{2.4\times10^{-7}}$ \\
1.5 & 96 & $\ensuremath{3.74\times10^{-73}}$ & $\ensuremath{1.31\times10^{-8}}$ \\
2 & 96 & $\ensuremath{5.72\times10^{-57}}$ & $\ensuremath{1.24\times10^{-11}}$ \\
3 & 96 & $\ensuremath{9.65\times10^{-42}}$ & $\ensuremath{1.18\times10^{-11}}$ \\
4 & 96 & $\ensuremath{1.57\times10^{-33}}$ & $\ensuremath{8.91\times10^{-13}}$ \\
6 & 192 & $\ensuremath{7.01\times10^{-58}}$ & $\ensuremath{6.73\times10^{-14}}$ \\
8 & 192 & $\ensuremath{8.7\times10^{-48}}$ & $\ensuremath{2.6\times10^{-14}}$ \\}
\newcommand{\ResidualThreshVal}{\ensuremath{1.0\times10^{-25}}}
\newcommand{\ResidualWorstVal}{\ensuremath{1.57\times10^{-33}}}
\newcommand{\RlocCounterVal}{0.0364613}
\newcommand{\RlocZeroResidualVal}{-5/12288}
\newcommand{\ScalarOrderVal}{16}
\newcommand{\ScreenedConstVal}{0.09765625}
\newcommand{\SelectedLowerVal}{4}
\newcommand{\SixtyFourDevVal}{0.45}
\newcommand{\SwaveCorrVal}{-0.014326825}
\newcommand{\SwaveFracVal}{30.7}
\newcommand{\TestChargeGapVal}{0.8300975079803747}
\newcommand{\TranscriptionCountVal}{60}
\newcommand{\TwentyOneDevVal}{2.05}
\newcommand{\TwoKVal}{2.472099569735163}
\newcommand{\VmcExtrapVal}{0.04995}
\newcommand{\VmcValuesVal}{$Z=2$: 0.0435, $Z=4$: 0.0467, $Z=8$: 0.0472, $Z=16$: 0.0497}
\newcommand{\VzeroVal}{1.023977169320997}
\newcommand{\WitnessVal}{0.0896598148781}
\newcommand{\WoneVal}{3.168529547708747982368425649780}
\newcommand{\WprimeOneVal}{-0.297254445136212166443653429304}

\newcommand{\acorr}{a_{\mathrm{corr}}}
\newcommand{\Rtail}{\mathcal{R}_{\mathrm{tail}}}
\newcommand{\Wsup}{W_{\mathrm{sup}}}
\newcommand{\Emin}{E_{\min}}
\newcommand{\st}[1]{\textsc{#1}}

\begin{document}

\title{Provenance of the atomic correlation constant $a_{\mathrm{corr}}$ in the substrate framework:\\
what is derived, what is conditional, and what remains open}

\author{Eliot Gable}
\email{10015996+egable@users.noreply.github.com}
\affiliation{Independent Researcher}

\date{\today}

\begin{abstract}
This is a provenance record for the correlation constant $\acorr$ as the substrate corpus treats it, so that every statement can be checked against a file, a function and a commit. In the corpus $\acorr=\epsilon_2-\eta_2$ is the high-nuclear-charge, same-centre, two-electron second-order correlation coefficient of the Coulomb comparison model. $\eta_2=\tfrac{9}{32}\ln\tfrac34-\tfrac{13}{432}$ is derived in closed form; $\epsilon_2$ is a published value, not a substrate result, so the number $\acorr=\AcorrDiagVal$ is a diagnostic. The corpus writes $\acorr$ in two parts in several ways: a contact block fixed by the coincidence limit plus a non-contact remainder $\Rtail$; an s-wave part (\SwaveFracVal{} percent) plus a higher-partial-wave part (\HigherlFracVal{} percent), of which the coincidence limit fixes only the large-$l$ tail; and a charge-tail (single-source response) part plus a two-centre correlation part. None of these splits determines $\acorr$ from the master Lagrangian. The test-charge kernel is not the equal-source interaction (exact contact ratio $4(\sqrt2-1)/3$). The two-centre part rests on the static two-source Born-Infeld energy. Its frozen superposition trial is closed in the sense of an exact linear differential equation with an exact algebraic-period germ, so it is computable by series continuation; the true static minimum is fenced but has no series; and the identification of either with the physical correlation of moving modes is a registered commitment that is not discharged. The physical $\acorr$ remains open. We list the routes tried and the reasons each does not close, with their sources.
\end{abstract}

\maketitle

\section{Purpose and scope}
The owner's summary of this line of work is that $\acorr$ ``was derived into two parts with $R_{\mathrm{tail}}$ closed as a series''. This paper sets out what the corpus actually establishes about that statement, with the evidence for each step, and reports where the statement aligns, partly aligns, or conflicts with the record. It is a provenance document: it introduces no new derivation, and every number is a macro generated by \path{papers/two_center_born_infeld/generate_numbers.py} from a corpus receipt at a recorded commit, or recomputed from a closed form. That generator also verifies that each of the \TranscriptionCountVal{} formula strings transcribed into the two companion papers occurs verbatim in its cited source, and that each cited function exists. The references are origin/dev at \RefDev, the frozen-ODE branch at \RefOde, the independent order-17 check at \RefCTwoSeven, the bond-research branch at \RefBond{} and the a\_corr lane branch at \RefAcorrTwoEight{} (full hashes in \path{papers/two_center_born_infeld/provenance_refs.json}). The static two-centre problem is treated in full in the companion paper~\cite{TwoCenterPaper}; here it appears only as one part of $\acorr$.

\subsection{Standing context}
The substrate framework is a Born-Infeld Lagrangian~\cite{BornInfeld1934} describing an elastic solid with an excited $E_8\times E_8$ dyad. $E_8\times E_8$ is observationally overdetermined, $D_{16}$ is excluded (\path{book/gr_qm_unification/sympy/substrate_theorems/rank16_observational_overdetermination.py}), and that includes its twist and spinor sector, its rank eight and its $D_8$ map. The framework has forward-derived quantization, entanglement and relativistic time dilation and is not classical. The particle mass equations are fixed closed forms in \path{book/scripts/substrate_mass_spectrum.py}. Particles are localized rotating wave modes of finite radius: two electrons in an atom are two lasting ripples of one medium, and their correlation is the discount in the energy of their joint, interfering pattern relative to the sum of the separate patterns. The Coulomb comparison model used below to define $\acorr$ is an imported reference problem, not the substrate's account of that discount.

\subsection{Status vocabulary}
\st{definition}; \st{derived} (forward from stated premises; the premises are named); \st{conditional} (rests on a registered commitment); \st{certified-exact}; \st{diagnostic} (a computation inside an imported comparison model or with literature inputs); \st{literature} (an external value or theorem); \st{open}; \st{withdrawn} (asserted earlier, later retracted in the corpus). Table~\ref{tab:claims} gives each claim's source file, function where one exists, and the last commit touching it; a dagger marks a commit not contained in origin/dev.

\section{What \texorpdfstring{$a_{\mathrm{corr}}$}{a_corr} is in the corpus}

\subsection{Definition}
For two electrons bound to one nucleus of charge $Z$ in the nonrelativistic Coulomb problem, the ground energy has the large-$Z$ expansion $E(Z)=-Z^2+\tfrac58Z+\epsilon_2+O(1/Z)$ in hartree, and the Hartree-Fock energy has the same first two terms and second-order coefficient $\eta_2$. The corpus defines
\begin{equation}
\acorr=\epsilon_2-\eta_2
\end{equation}
(\path{atomic_correlation_constant_not_forced.py::verdict}; \st{definition}). The $1/Z$ expansion and the Hartree-Fock reference are perturbative constructs of the imported Coulomb model; in the substrate framework they serve only as a comparison problem, under the quarantine of perturbation theory as a method, and a substrate-native correlation must come from the moving-mode energy functional of Sec.~\ref{sec:rtail} and the open items. It is a one-centre constant: it carries no dependence on a separation between centres, so it is not the two-centre correlation energy that a bond depth needs. The $O(Z)$ terms are equal because the exchange-lowered spatially antisymmetric partner of $1s^2$ is excluded in both curves (\path{pauli_shared_oz_term}), and the static monopole content $J=\tfrac58Z$ cancels in the difference.

\subsection{The closed part: \texorpdfstring{$\eta_2$}{eta_2}}
The first-order orbital equation $(-\tfrac12\nabla^2-1/\rho+\tfrac12)\phi_1=-(V_C-\tfrac58)\phi_0$ with $V_C(r)=(1-(1+r)e^{-2r})/r$ is solved by $\phi_1=f\phi_0$ in the Dalgarno-Lewis reduction~\cite{DalgarnoLewis1955}, which turns it into a first-order equation for $w=f'$ whose solution regular at the origin is elementary. The energy integral then gives
\begin{equation}
\eta_2=\frac{9}{32}\ln\frac34-\frac{13}{432}=\EtaTwoVal ,
\end{equation}
the logarithm arising by Frullani integration of the $1/r$ part of $f'$ (\path{eta2_symbolic_derivation}; establish-or-bail checks of the defining equation, regularity, decay and the self-adjoint identity). It agrees with the published closed form of Loos and Gill~\cite{LoosGill2009}, which is an oracle here. Status: \st{derived}, inside the imported Coulomb operator, whose kinetic and potential terms are not derived from the Born-Infeld medium in this calculation.

\subsection{The open part: \texorpdfstring{$\epsilon_2$}{epsilon_2}}
The corpus value $\epsilon_2=\EpsTwoVal$ is read from the publication of Wang~\cite{Wang2013} (\path{second_order_partial_waves.py::published_eps2}; \st{literature}). Hence the corpus number
\begin{align}
\acorr&=\EpsTwoVal-(\EtaTwoVal)\nonumber\\
&=\AcorrDiagVal ,\qquad |\acorr|=\AcorrMagVal ,
\end{align}
is a \st{diagnostic} that inherits the provenance of the published $\epsilon_2$; it is not a substrate prediction. The corpus reduces $\epsilon_2$ to a sum over total-$L=0$ partial-wave channels. The s-wave particular solution $\tfrac58r_>-\tfrac38\ln r_>+\tfrac{3}{16r_>}$ is exact; the large-$l$ tail follows the Schwartz law $-\tfrac{45}{256}(l+\tfrac12)^{-4}+\tfrac{225}{1024}(l+\tfrac12)^{-6}+\dots$~\cite{Schwartz1962}, whose sum is elementary through half-integer Hurwitz values; the matching data at the crossover shell $r_1=r_2$ are fixed in closed form by a field junction. The crossover correction, however, is not in the five-dimensional elementary-plus-exponential-integral homogeneous space of the channel operator, so the s-wave channel function is not of that elementary form and $\epsilon_2(0)$ does not close by this assembly; numerically $\epsilon_2(0)=-0.125334(1)$ (\path{eps2_swave_crossover_is_non_elementary}). This is a non-membership result for one assembly, not an impossibility proof for a closed form. Status of a closed form for $\epsilon_2$: \st{open}.

\section{The two-part decompositions}
The corpus writes $\acorr$ as two parts in three different senses. They must not be conflated, and the symbol $R_{\mathrm{tail}}$ carries different meanings across them (written $\Rtail$ here for the remainder of the first split).

\subsection{Contact block plus non-contact remainder}
Write the target as a global integral $\acorr=\int_0^\infty G(r_{12})\,w(r_{12})\,dr_{12}$ of a two-centre field-energy density against a pair measure, choose an overlap radius $r_*$, and split
\begin{equation}
\acorr=\sum_k\mu_kG_k+\Rtail,\qquad \mu_k=\int_0^{r_*}r^{p_k}w\,dr,\qquad \Rtail=\int_{r_*}^{\infty}G\,w\,dr,
\end{equation}
where $G(r)=\sum_kG_kr^{p_k}$ near coalescence (\path{acorr_proximity_cap_rank_analysis.py}). Matching the global density to the local coincidence hole, $G_k=L_k$, is a square system: one equation and one open contact coefficient per order, with no surplus equation, so the match transfers local moments into contact data and places no constraint on the global integral (\path{substrate_correlation_coincidence_match_underdetermines}; \st{derived}). A Born-Infeld proximity cap at a minimum separation is a contact-local boundary condition: symbolically and over 200 random instances its row lies in the span of the match rows, its projection on the $(\acorr,\Rtail)$ columns vanishes, and the $(\acorr,\Rtail)$ subsystem stays rank one in two unknowns; the cap's consistency fixes the minimum separation, not $\acorr$ (\path{substrate_correlation_proximity_cap_is_contact_local_no_rank}; \st{derived}). This split is a bookkeeping identity; it shows that no local condition reaches $\acorr$, and it names $\Rtail$ as the residual carried by the global two-centre field.

\subsection{s-wave plus higher partial waves}
With literature per-$l$ increments of the Coulomb problem (Carroll, Silverstone and Metzger~\cite{CSM1979}) the corpus finds the s-wave correlation $\epsilon_2(0)-\eta_2=\SwaveCorrVal$, \SwaveFracVal{} percent of $\acorr$, and \HigherlFracVal{} percent in $l\ge1$ (\path{acorr_partialwave_decomposition.py}; \st{diagnostic}). The coincidence limit $r_{12}\to0$ is the large-$l$ limit: it reproduces the shared cusp law, whose sum from $l=0$ is $+\CuspTailSumVal$, nowhere near $\epsilon_2$, while most of $\acorr$ sits in the low-$l$ head that the coincidence data never reach. The resisting content is the short-range Kato coalescence cusp~\cite{Kato1957}, whose branch point gives the partial-wave expansion its power-law tail.

\subsection{Charge-tail part plus two-centre correlation}
\paragraph{Charge tail.} The single-source Born-Infeld potential and the linear response of an infinitesimal test charge are evaluated in \path{acorr_charge_tail_dependency_FINDINGS.md}. For $s\ge0$ and core scale $r_0>0$, the substitutions $r=r_0(1-t)/t$, $t=au$, $a=r_0/(s+r_0)$ give
\begin{equation}
U(s)=\frac{1}{s+r_0}\int_0^1\frac{du}{\sqrt{(1-au)^4+(au)^4}},\qquad
\frac1s-\frac{r_0^4}{10s^5}\le U(s)\le\frac1s ,
\end{equation}
with the quartic in $[1/8,1]$ on the fixed interval, which removes a scale-dependent quadrature failure (\st{derived}, scalar route only; landed numerical-domain component). The same record states the scope plainly: the mechanical displacement, strain, scalar potential and conjugate flux are distinct types, a scalar $1/r$ potential does not identify a compatible mechanical strain, and no physical $\acorr$ follows. The decisive check is at contact. By displacement-preserving dilation a source of charge $q$ has energy $Kq^{3/2}$ with $K=\Gamma(1/4)^2/(6\sqrt\pi)=\KVal$, so the interaction of two finite sources at contact is $V(q_1,q_2;0)=K[(q_1+q_2)^{3/2}-q_1^{3/2}-q_2^{3/2}]$, while the test-charge coefficient is its one-sided derivative. For two unit sources
\begin{equation}
\frac{V(1,1;0)}{U(0)}=\frac{4(\sqrt2-1)}{3}=\EndpointRatioVal\neq1,\qquad U(0)-V(1,1;0)=\Bigl(\frac72-2\sqrt2\Bigr)K=\TestChargeGapVal ,
\end{equation}
(\path{bi_test_charge_pair_energy.py::contact_interaction_certificate}; \st{certified-exact}). The isolated-source test-charge potential is therefore not the equal-source interaction already in the static point-source diagnostic, and the partial-wave calculations that used it as a pair kernel are comparison calculations whose transfer to two finite sources is not established. The static linear far-field strain of a spherical source is pure longitudinal and pure monopole; inside the Coulomb comparison operator its content is the $l=0$ direct term that cancels in $\acorr$, and the retarded magnetic wake enters at $O(\alpha^2)\sim5\times10^{-5}$ (\path{substrate_correlation_strain_tail_and_wake_not_the_resisting_piece}; \st{diagnostic}). This does not exclude wake, twist or retardation contributions to the physical energy of finite moving modes; the framework holds that those effects are always present, and their size must be computed, not assumed.

\paragraph{Two-centre correlation.} The second part is the energy of the joint two-source field, which is where Born-Infeld nonsuperposition acts. The corpus registers its identification with $\acorr$ as the commitment \path{rtail_two_center_m11_physical_identification}: mapping the static Euclidean two-centre energy to the correlation of moving atomic modes requires the finite moving-mode reduction (observer clock and measure, retarded wake, angular momentum and twist, interference), which is not built (\st{conditional}). The M11 route validation, which tried to compute $\acorr$ from the true two-electron field minimizer, returns open with that blocking antecedent named, \texttt{full\_elastic\_observer\_energy is None}, together with the missing physical two-electron bound mode; it also records that softening the interaction alone, with the nonrelativistic framework fixed, shifts $\acorr$ by $O(r_0^2)\sim10^{-9}$ at atomic scale, so the leading part reduces to the Coulomb $\epsilon_2-\eta_2$ unless the moving-mode terms change it (\path{acorr_m11_true_minimizer_route_validation.py}; \st{open}).

\section{\texorpdfstring{$R_{\mathrm{tail}}$}{R_tail} and its series}
\label{sec:rtail}
In the files named \path{rtail_*} the symbol $R_{\mathrm{tail}}(d)$ denotes the static two-source Born-Infeld energy at separation $d$, either of the true minimizer $\Emin(d)$ or, in older files, of the frozen superposition trial $\Wsup(d)$ (the reduction once labelled the true $R_{\mathrm{tail}}$ was relabelled as the frozen trial; see the timeline). The companion paper~\cite{TwoCenterPaper} gives the full account; the results that bear on the owner's statement are these.
\begin{enumerate}
\item The frozen trial $\Wsup$ satisfies an exact linear differential equation of order \EnergyOrderVal{} with integer polynomial coefficients, certified by exact divergence identities (replayed for this record), and a smaller order-\OrderSeventeenVal{} annihilator has independent source clearance on its branch. The physical branch has order at least \SelectedLowerVal{} and is not algebraic. An exact germ of eighteen algebraic periods at $d=1$ selects the solution, so $\Wsup$ is computable to arbitrary precision by series continuation once the germ is enclosed to that precision; a recorded continuation step agrees with independent quadrature to \IvpErrVal{}, and an unreviewed multi-step continuation to \BondDiffMaxVal{} or better on $1.25\le d\le4$ (observed agreements, not error bars). In this precise sense the frozen trial is ``closed as a series''.
\item The frozen trial is a strict upper bound on the static energy, never the energy itself. For the true minimum the corpus has certified fences (the contact force $1/2$ from a contact enclosure, the interval $\DoneLowerVal\le V(1)\le\DoneUpperVal$, a dual lower bound) but no differential equation and no series; its large-separation law is conditional on a registered commitment.
\item Neither static object has been identified with the physical two-mode correlation (the commitment above), and neither is the remainder $\Rtail$ of the first split, which is an integral of a density that the corpus has not constructed.
\end{enumerate}
The former claim that the physical two-centre energy is a proved non-elementary order-four Picard-Fuchs period is \st{withdrawn}: \path{rtail_known_problems_context.md} states that the former claim is withdrawn and that the checked algebra concerns the frozen trial. The derivations README at the recorded origin/dev commit still carried the older ``now PROVEN NON-ELEMENTARY'' paragraph; the corpus commit that accompanies this paper replaces it with the current status, and corrects the two-centre synthesis module, which had called the true large-separation law exact.

\section{Routes tried and their outcomes}
\begin{longtable}{lll}
\caption{Routes to $\acorr$ recorded in the corpus.}\label{tab:routes}\\
\toprule
\cell{0.27}{route} & \cell{0.42}{outcome} & \cell{0.24}{source}\\
\midrule
\endfirsthead
\toprule
\cell{0.27}{route} & \cell{0.42}{outcome} & \cell{0.24}{source}\\
\midrule
\endhead
\cell{0.27}{Constant installed in the forward element generator (August)} & \cell{0.42}{the original 2p diagnostic installs $0.04667$ hartree as an input; its forward provenance was not established, and a later commit separated it from the cross-shell correlation $E_c$} & \cell{0.24}{a\_corr lane journal; timeline}\\
\cell{0.27}{$\epsilon_2$ by Green's-function assembly} & \cell{0.42}{crossover correction outside the elementary plus Ei span; $\epsilon_2(0)$ does not close by this assembly (non-membership)} & \cell{0.24}{\path{eps2_swave_crossover_is_non_elementary}}\\
\cell{0.27}{Rigid limit versus deformed metric} & \cell{0.42}{the lapse rationalizes the source but the acoustic-measure root cancels it; the universal cancellation was later withdrawn and the observer measure must be derived} & \cell{0.24}{\path{substrate_correlation_hole_is_a_deformed_metric_problem}}\\
\cell{0.27}{Global field-energy route} & \cell{0.42}{the saturation root resums globally and the single-centre building blocks are $\Gamma(1/4)$-clean; the correlation functional and its field-energy map remain open} & \cell{0.24}{\path{substrate_correlation_global_field_energy_route}}\\
\cell{0.27}{Coincidence match} & \cell{0.42}{square system; under-determines $\acorr$} & \cell{0.24}{\path{substrate_correlation_coincidence_match_underdetermines}}\\
\cell{0.27}{Static strain tail and wake} & \cell{0.42}{inside the comparison operator: the static tail is $l=0$ and cancels, the wake is $O(\alpha^2)$; not a statement about finite moving modes} & \cell{0.24}{\path{substrate_correlation_strain_tail_and_wake_not_the_resisting_piece}}\\
\cell{0.27}{Cuspless test-charge kernel} & \cell{0.42}{comparison calculation; the kernel is not the equal-source interaction} & \cell{0.24}{\path{substrate_native_cuspless_correlation_constant}}\\
\cell{0.27}{Proximity cap} & \cell{0.42}{contact-local, adds no rank} & \cell{0.24}{\path{substrate_correlation_proximity_cap_is_contact_local_no_rank}}\\
\cell{0.27}{Order-four Picard-Fuchs closure of $R_{\mathrm{tail}}$} & \cell{0.42}{withdrawn; downgraded to a rank-six upper bound with open obligations, then superseded by the exact order-18 equation for the frozen trial} & \cell{0.24}{\path{rtail_known_problems_context.md}}\\
\cell{0.27}{M11 true-minimizer route} & \cell{0.42}{open; the moving-mode two-electron functional and the physical bound mode are missing} & \cell{0.24}{\path{acorr_m11_true_minimizer_route_validation.py}}\\
\cell{0.27}{Harmonium pair law} & \cell{0.42}{not a forced universal constant} & \cell{0.24}{\path{harmonium_correlation_not_forced.py}}\\
\bottomrule
\end{longtable}

The module title \path{atomic_correlation_constant_not_forced.py} records the cumulative result: its verdict exports \texttt{"a\_corr": None} and \texttt{"physical\_generator\_complete": False}. ``Not forced'' there means that no route recorded so far forces a value; it is not a proof that no route can.

\section{What is certified, and to what precision}
\begin{itemize}
\item Exact: $\eta_2$ in closed form; the contact ratio $4(\sqrt2-1)/3$; the endpoint energies $2K$ and $2^{3/2}K$ of the static two-source problem; the order-18 annihilator of the frozen trial and its indicial data; the frozen large-separation law $\Wsup=2K+1/d-(3K/4)/d^4+o(d^{-4})$; the lower bound \SelectedLowerVal{} on the frozen trial's order and its nonalgebraicity.
\item Interval-certified: the true static minimum's contact enclosure on $0\le d\le1/2$ and its interval at $d=1$.
\item Observed (agreement of routes, not error bars): the frozen trial's values and continuation agreements; the numerical $\epsilon_2(0)$.
\item Literature inputs: $\epsilon_2$ and the per-$l$ increments used in the partial-wave split.
\end{itemize}

\section{Assessment of the owner's statement}
The statement partly aligns with the record. It aligns in that the corpus does split $\acorr$ into two parts, in each of the three senses above, and in that the frozen two-centre trial energy is closed in the sense of an exact differential equation with an exact germ, computable to arbitrary precision by series continuation. It conflicts in three places. First, none of the splits is a derivation of $\acorr$'s value: the contact-plus-remainder split is a bookkeeping identity whose remainder is unconstructed, the partial-wave split uses published Coulomb values, and the charge-tail-plus-two-centre split rests on an identification that is a registered, undischarged commitment. Second, what is closed as a series is the frozen superposition trial, an upper bound, not the true static minimum and not the correlation of moving modes. Third, the corpus value of $\acorr$ inherits the published $\epsilon_2$; it is a diagnostic, not a substrate result. The physical $\acorr$ is open.

\section{What remains open}
\begin{enumerate}
\item The physical moving-mode two-electron correlation functional, with its observer clock and measure (by the owner's ruling, one global tick and the actual local voxel volume), retarded strain field, motion wake, twist and interference, and its identification with $\acorr$; and the physical bound mode it requires.
\item The charge sector of the master Lagrangian's cross term (phase-carrying mode class, source coupling, $1/r$ tail and Maxwell regime), which the a\_corr lane orders before the bound mode.
\item A closed form, or a proof of impossibility, for $\epsilon_2$ and for $\epsilon_2(0)$.
\item The remainder $\Rtail$ as an actual integral of a constructed density.
\item For the static two-centre energy: the items listed in~\cite{TwoCenterPaper}, including the minimal order of the frozen trial and any series or equation for the true minimum.
\item Verification of corpus citation fields not checked for this record, for example the forcing node's statement of what Carley and Kiessling establish about the linear Bopp-Podolsky potential.
\end{enumerate}

\appendix
\section{Claims, status and sources}
Paths are relative to \path{book/gr_qm_unification/sympy/}.
\begingroup\small
\begin{longtable}{lllll}
\caption{Claims of this paper with status and source.}\label{tab:claims}\\
\toprule
\cell{0.04}{id} & \cell{0.33}{claim} & \cell{0.17}{status} & \cell{0.29}{source} & \cell{0.1}{commit}\\
\midrule
\endfirsthead
\toprule
\cell{0.04}{id} & \cell{0.33}{claim} & \cell{0.17}{status} & \cell{0.29}{source} & \cell{0.1}{commit}\\
\midrule
\endhead
\cell{0.04}{A1} & \cell{0.33}{a\_corr is defined as eps\_2 - eta\_2 of the Coulomb comparison model (high-Z, same centre, two electrons)} & \cell{0.17}{DEFINITION} & \cell{0.29}{\path{substrate_theorems/atomic_correlation_constant_not_forced.py::verdict}} & \cell{0.1}{\texttt{56d7c16cfe}} \\
\cell{0.04}{A2} & \cell{0.33}{eta\_2 = (9/32)ln(3/4) - 13/432 by the Dalgarno-Lewis reduction} & \cell{0.17}{DERIVED (imported Coulomb operator)} & \cell{0.29}{\path{substrate_theorems/atomic_correlation_constant_not_forced.py::eta2_symbolic_derivation}} & \cell{0.1}{\texttt{56d7c16cfe}} \\
\cell{0.04}{A3} & \cell{0.33}{eps\_2 = published value (Wang 2013), not a substrate result} & \cell{0.17}{LITERATURE} & \cell{0.29}{\path{imported_theorems/second_order_partial_waves.py::published_eps2}} & \cell{0.1}{\texttt{56d7c16cfe}} \\
\cell{0.04}{A4} & \cell{0.33}{s-wave and l$>$=1 split of a\_corr from literature per-l increments} & \cell{0.17}{DIAGNOSTIC (literature inputs)} & \cell{0.29}{\path{derivations/acorr_partialwave_decomposition.py}} & \cell{0.1}{\texttt{674a28b5fd}} \\
\cell{0.04}{A5} & \cell{0.33}{the s-wave crossover correction lies outside the elementary plus Ei span (non-membership only)} & \cell{0.17}{DERIVED (conditional)} & \cell{0.29}{\path{substrate_theorems/atomic_correlation_constant_not_forced.py::eps2_swave_crossover_is_non_elementary}} & \cell{0.1}{\texttt{56d7c16cfe}} \\
\cell{0.04}{A6} & \cell{0.33}{coincidence match is a square system; it fixes contact data, not a\_corr} & \cell{0.17}{DERIVED} & \cell{0.29}{\path{substrate_theorems/atomic_correlation_constant_not_forced.py::substrate_correlation_coincidence_match_underdetermines}} & \cell{0.1}{\texttt{56d7c16cfe}} \\
\cell{0.04}{A7} & \cell{0.33}{proximity cap adds zero rank; (a\_corr, R\_tail) block stays rank 1 in 2 unknowns} & \cell{0.17}{DERIVED} & \cell{0.29}{\path{derivations/acorr_proximity_cap_rank_analysis.py}} & \cell{0.1}{\texttt{56d7c16cfe}} \\
\cell{0.04}{A8} & \cell{0.33}{test-charge contact ratio V(1,1;0)/U(0) = 4(sqrt2-1)/3, not 1} & \cell{0.17}{CERTIFIED-EXACT} & \cell{0.29}{\path{substrate_theorems/bi_test_charge_pair_energy.py::contact_interaction_certificate}} & \cell{0.1}{\texttt{b436d403d6}} \\
\cell{0.04}{A9} & \cell{0.33}{scalar point-source potential: exact substitution, quartic range [1/8,1], Coulomb bracket} & \cell{0.17}{DERIVED (scalar route only)} & \cell{0.29}{\path{derivations/acorr_charge_tail_dependency_FINDINGS.md}} & \cell{0.1}{\texttt{cfa43b6430}} \\
\cell{0.04}{A10} & \cell{0.33}{static linear strain tail is l=0 and the wake is O(alpha\^{}2) inside the Coulomb comparison only} & \cell{0.17}{CONDITIONAL DIAGNOSTIC} & \cell{0.29}{\path{substrate_theorems/atomic_correlation_constant_not_forced.py::substrate_correlation_strain_tail_and_wake_not_the_resisting_piece}} & \cell{0.1}{\texttt{56d7c16cfe}} \\
\cell{0.04}{A11} & \cell{0.33}{static two-centre energy to a\_corr identification} & \cell{0.17}{COMMITMENT (M11-gated)} & \cell{0.29}{\path{substrate_theorems/rtail_two_center_synthesis.py::two_center_synthesis}} & \cell{0.1}{\texttt{56d7c16cfe}} \\
\cell{0.04}{A12} & \cell{0.33}{M11 true-minimizer route: physical moving-mode functional and bound mode missing} & \cell{0.17}{OPEN} & \cell{0.29}{\path{derivations/acorr_m11_true_minimizer_route_validation.py}} & \cell{0.1}{\texttt{56d7c16cfe}} \\
\cell{0.04}{A13} & \cell{0.33}{former order-4 Picard-Fuchs non-elementarity claim} & \cell{0.17}{WITHDRAWN} & \cell{0.29}{\path{derivations/rtail_known_problems_context.md}} & \cell{0.1}{\texttt{745293d9c5}} \\
\cell{0.04}{A14} & \cell{0.33}{physical a\_corr value from the master Lagrangian} & \cell{0.17}{OPEN} & \cell{0.29}{\path{substrate_theorems/atomic_correlation_constant_not_forced.py::verdict}} & \cell{0.1}{\texttt{56d7c16cfe}} \\
\bottomrule

\end{longtable}
\endgroup

\section{Dated history}
Commits are found by exact subject in the remote history; times are UTC; a dagger marks a commit not contained in origin/dev.
\begingroup\small
\begin{longtable}{lll}
\caption{The history of $\acorr$ and $R_{\mathrm{tail}}$ in the corpus.}\label{tab:timeline}\\
\toprule
\cell{0.17}{UTC} & \cell{0.14}{commit} & \cell{0.62}{subject}\\
\midrule
\endfirsthead
\toprule
\cell{0.17}{UTC} & \cell{0.14}{commit} & \cell{0.62}{subject}\\
\midrule
\endhead
\cell{0.17}{2026-08-11 18:23Z} & \cell{0.14}{\texttt{74bebdae80}} & \cell{0.62}{feat(book): the COMPLETE forward generator ; zero fitted constants, 2p row $<$1\%} \\
\cell{0.17}{2026-08-11 21:27Z} & \cell{0.14}{\texttt{a4dc3e19eb}} & \cell{0.62}{fix: correct the E\_c=a\_corr over-claim ; verified they are different quantities} \\
\cell{0.17}{2026-09-17 02:30Z} & \cell{0.14}{\texttt{b2a7b01ba6}} & \cell{0.62}{feat(gr\_qm): a\_corr is measured, not derivable -{}- and the corpus disagrees with itself by 2} \\
\cell{0.17}{2026-09-18 01:32Z} & \cell{0.14}{\texttt{1ae99b7475}} & \cell{0.62}{feat(gr-qm): eps\_2(0) does not close by the Green's-function assembly -{}- the crossover correction is non-elementary} \\
\cell{0.17}{2026-09-18 05:31Z} & \cell{0.14}{\texttt{5fbbc005ff}} & \cell{0.62}{research(gr-qm): prove the rigid-limit source non-elementary; the substrate answer is the deformed metric} \\
\cell{0.17}{2026-09-18 05:59Z} & \cell{0.14}{\texttt{0fc13294cf}} & \cell{0.62}{research(gr-qm): the substrate a\_corr is a global field energy; the saturation root resums globally} \\
\cell{0.17}{2026-09-18 07:08Z} & \cell{0.14}{\texttt{b99c392aa1}} & \cell{0.62}{research(gr-qm): name the two-center Born-Infeld field energy as the open a\_corr residual} \\
\cell{0.17}{2026-09-18 14:50Z} & \cell{0.14}{\texttt{674a28b5fd}} & \cell{0.62}{research(sympy): land a\_corr coincidence match derivations and two open verdicts} \\
\cell{0.17}{2026-09-18 15:58Z} & \cell{0.14}{\texttt{caf6916754}} & \cell{0.62}{research(sympy): land proximity-cap contact-local no-rank verdict for a\_corr} \\
\cell{0.17}{2026-09-18 17:19Z} & \cell{0.14}{\texttt{4626c0041a}} & \cell{0.62}{research: land two-center Born-Infeld field energy R\_tail closed-form bracket} \\
\cell{0.17}{2026-09-18 20:06Z} & \cell{0.14}{\texttt{a2ea24ca21}} & \cell{0.62}{research: prove the two-center Born-Infeld field energy R\_tail is non-elementary} \\
\cell{0.17}{2026-09-19 01:27Z} & \cell{0.14}{\texttt{6744384080}} & \cell{0.62}{research: downgrade R\_tail two-center BI structure to a rank-6 upper bound with open obligations} \\
\cell{0.17}{2026-09-19 03:24Z} & \cell{0.14}{\texttt{745293d9c5}} & \cell{0.62}{fix(physics): derive R\_tail trial bounds and correct correlation claims} \\
\cell{0.17}{2026-09-21 05:09Z} & \cell{0.14}{\texttt{409d83bbff}} & \cell{0.62}{feat(gr-qm): two-center Born-Infeld minimizer is unique and two-sided bracketed} \\
\cell{0.17}{2026-09-21 14:34Z} & \cell{0.14}{\texttt{35a424b90d}} & \cell{0.62}{docs(gr-qm): label the lambda-reduction as the frozen superposition energy, not the true R\_tail} \\
\cell{0.17}{2026-09-21 15:10Z} & \cell{0.14}{\texttt{826a804f28}} & \cell{0.62}{feat(gr-qm): validate the M11 true-minimizer route to a\_corr - OPEN, blocking antecedent named} \\
\cell{0.17}{2026-10-28 23:16Z} & \cell{0.14}{\texttt{33b095be04}$^{\dagger}$} & \cell{0.62}{docs: record acorr derivation history and remaining antecedents} \\
\cell{0.17}{2026-11-01 03:55Z} & \cell{0.14}{\texttt{ef75344b6c}} & \cell{0.62}{feat(gr-qm): establish selected physical energy order lower bound} \\
\cell{0.17}{2026-11-01 12:06Z} & \cell{0.14}{\texttt{e3b4470853}$^{\dagger}$} & \cell{0.62}{feat: derive order17 frozen energy annihilator} \\
\cell{0.17}{2026-11-01 13:13Z} & \cell{0.14}{\texttt{9188e5a275}$^{\dagger}$} & \cell{0.62}{feat: classify order17 apparent divisors} \\
\cell{0.17}{2026-11-01 14:03Z} & \cell{0.14}{\texttt{703d37e793}$^{\dagger}$} & \cell{0.62}{test: independently check order17 divisor and endpoint data} \\
\bottomrule

\end{longtable}
\endgroup

\section{Premise ledger}
Consumed by the definition and the diagnostic numbers: the nonrelativistic Coulomb comparison problem and its $1/Z$ expansion (imported reference model); the published $\epsilon_2$ and per-$l$ increments (literature). Consumed by the static two-centre results: the Born-Infeld energy density (the chain's axiom), unit point sources in static Euclidean space (definitions of the diagnostic) and the finite-energy admissible class (commitment). Registered commitments that bound the conclusions: the M11 identification of the static two-centre energy with $\acorr$, and, for the true large-separation law, the induced-dipole structure. The $E_8\times E_8$ geometry, its $D_8$ placement and the particle equations are retained in the framework and are not consumed by these calculations.

\bibliography{acorr_derivation}

\end{document}
